\chapter{Innovation Project, Re-Design and Product Presentation}
\label{chap:innovation_redesign_presentation}

\begin{tcolorbox}[enhanced, colback=boxnavybg, colframe=brandnavy, boxrule=1.2pt, arc=4pt, title={\Large\bfseries\sffamily Unit 6 Overview \& Learning Objectives}]
\sffamily\normalsize
This unit synthesizes the complete Design Thinking lifecycle into iterative refinement, structured re-design, and high-impact stakeholder communication. By the end of this chapter, students will be able to:
\begin{itemize}[leftmargin=1.2em, itemsep=2pt, topsep=3pt]
    \item Execute the \textbf{Design Iteration Loop} from feedback analysis to re-testing.
    \item Contrast \textbf{Design Refactoring} (improving internal structure) with \textbf{Product Re-Design} (overhauling core architecture).
    \item Apply the feedback priority formula: $P = F \times S \times C$.
    \item Recognize triggers for \textbf{Strategic Pivoting} vs. Continuous Improvement using the \textbf{PDCA Cycle}.
    \item Tailor product presentations to specific audiences: users, engineers, and executive investors.
    \item Structure value propositions using SRI International's \textbf{NABC Model} (Need, Approach, Benefits per costs, Competition).
    \item Deliver compelling 30--60 second \textbf{Elevator Pitches}.
    \item Assemble a complete \textbf{Technical Design Case Study} with requirement traceability.
\end{itemize}
\end{tcolorbox}

\section{Orientation: The Iterative Refinement Lifecycle}

Innovation in Design Thinking is an ongoing cycle of empirical refinement. A product is never truly complete after the first build; it matures through disciplined feedback analysis, refactoring, and clear stakeholder advocacy.

\begin{keyequation}[The Iteration Progression Cycle]
\begin{equation}
\text{Define Assumption} \longrightarrow \text{Build Prototype} \longrightarrow \text{Test with Users} \longrightarrow \text{Analyze Evidence} \longrightarrow \text{Refactor / Re-design} \longrightarrow \text{Retest}
\end{equation}
\end{keyequation}

\section{Iteration, Refactoring, and Strategic Pivoting}

\subsection{Design Refactoring vs. Product Re-Design}
\begin{itemize}[leftmargin=1.2em, itemsep=3pt]
    \item \textbf{Design Refactoring:} Improving the internal consistency, layout alignment, typography hierarchy, component modularity, or code structure of a product \textbf{without altering its fundamental functional behavior}.
    \item \textbf{Product Re-Design:} Fundamental structural and functional overhaul triggered when user testing proves that the foundational problem definition, value proposition, or mechanical/software architecture has failed.
    \item \textbf{Regression Testing:} After any refactoring or re-design, the team must retest previously validated workflows to ensure new modifications did not break existing features.
\end{itemize}

\subsection{User Feedback Prioritization Formula}
To prevent subjective debates about which user critiques to implement first, engineering teams calculate a quantitative Priority Score:

\begin{keyequation}[Priority Score Calculation]
\begin{equation}
P = F \times S \times C
\end{equation}
\begin{itemize}[leftmargin=1.2em, itemsep=2pt]
    \item $F \in [1, 5]$: **Frequency** of the issue observed across test participants.
    \item $S \in [1, 5]$: **Severity** of the issue ($1 = \text{minor cosmetic preference}$; $5 = \text{complete task blocker / safety hazard}$).
    \item $C \in [0.1, 1.0]$: **Confidence** in the research evidence.
\end{itemize}
\end{keyequation}

\subsection{Strategic Pivots vs. Continuous Improvement (PDCA)}
\begin{itemize}[leftmargin=1.2em, itemsep=3pt]
    \item \textbf{Strategic Pivot:} A structured course correction testing a new fundamental hypothesis regarding customer segment, delivery channel, or value proposition while preserving validated past learnings (e.g., YouTube pivoting from video dating to general video sharing).
    \item \textbf{Continuous Improvement (PDCA Cycle):} Incremental quality enhancement following:
    $$\textbf{Plan} \;\rightarrow\; \textbf{Do} \;\rightarrow\; \textbf{Check} \;\rightarrow\; \textbf{Act}$$
\end{itemize}

\section{Product Presentation \& Design Communication}

\subsection{Audience-Centred Communication}
An effective pitch is calibrated specifically to the stakeholder's primary decision criteria:
\begin{table}[h!]
\centering
\small
\renewcommand{\arraystretch}{1.3}
\begin{tabularx}{\linewidth}{p{3.2cm}X X}
\toprule
\textbf{Audience} & \textbf{Primary Concerns} & \textbf{Presentation Focus} \\
\midrule
\textbf{End Users} & Usability, convenience, personal value, privacy. & Intuitive workflows, comfort, immediate problem relief. \\
\textbf{Engineers / Builders} & Technical feasibility, scalability, safety, stack. & System architecture diagrams, component specs, latency, test metrics. \\
\textbf{Investors / Executives} & Business viability, market size, unit economics, ROI. & Customer acquisition cost, scalable revenue, competitive moat, NABC. \\
\bottomrule
\end{tabularx}
\caption{Tailoring Design Presentations to Key Stakeholders.}
\end{table}

\subsection{High-Impact Presentation Rules}
\begin{enumerate}[leftmargin=1.2em, itemsep=3pt]
    \item \textbf{The Compelling Hook:} Open with a startling real-world statistic, video clip, or user story that immediately demonstrates an urgent, unsolved human pain point.
    \item \textbf{The Rule of One:} Every slide must convey exactly ONE dominant concept with high-contrast, uncluttered visuals.
    \item \textbf{Live Representative Task Demo:} Demonstrate a complete, authentic user flow rather than clicking through boring settings menus (*``Show, Don't Tell''*).
    \item \textbf{Radical Candor (Stating Limitations):} Explicitly state known product limitations and future milestones. Acknowledging boundaries builds immense credibility with technical evaluators.
\end{enumerate}

\section{The NABC Pitching Model}

Developed by Curtis Carlson at SRI International, the NABC framework organizes innovation proposals into 4 fundamental pillars:
\begin{itemize}[leftmargin=1.2em, itemsep=3pt]
    \item \textbf{Need (N):} What is the specific, quantitative, validated customer pain point? Back it with user research evidence.
    \item \textbf{Approach (A):} What is your unique, feasible solution, technology mechanism, or service model?
    \item \textbf{Benefits per Costs (B):} What measurable benefits (time saved, errors reduced, revenue gained) does the user receive relative to the financial, behavioral, and time costs incurred?
    \item \textbf{Competition / Alternatives (C):} How does your solution provide clear relative advantage compared to existing market competitors, manual workarounds, and non-consumption?
\end{itemize}

\section{The Elevator Pitch}

\begin{definition}[Elevator Pitch]
A high-density 30--60 second spoken narrative designed to create sufficient clarity and excitement to secure a follow-up demonstration or partnership.
\end{definition}

\begin{keyequation}[The Universal Product Positioning Template]
\begin{quote}
\textit{``For \textbf{[target user]} who experiences \textbf{[validated problem]}, our \textbf{[product name]} is a \textbf{[product category]} that delivers \textbf{[compelling core benefit]}. Unlike \textbf{[primary competitor / current workaround]}, our solution \textbf{[key differentiator]}. User tests demonstrate \textbf{[validated metric / empirical evidence]}.''}
\end{quote}
\end{keyequation}

\section{The Technical Case Study Deliverable}

A comprehensive INT335 capstone design portfolio integrates the entire engineering lifecycle:
\begin{enumerate}[leftmargin=1.2em, itemsep=2pt]
    \item \textbf{Executive Summary \& Context}
    \item \textbf{Empathy Research Findings \& POV Definition}
    \item \textbf{Divergent Concepts \& Screening Matrix}
    \item \textbf{Prototype Evolution (Lo-Fi to Hi-Fi)}
    \item \textbf{Usability Testing Logs \& Feedback Capture Matrices}
    \item \textbf{Requirement Traceability Matrix}
    \item \textbf{Final Technical Architecture, Safety Specs \& Viability Plan}
\end{enumerate}

\section{Unit 6 Self-Assessment Test}

\begin{chaptertest}[Innovation Project, Re-Design \& Product Presentation]
\textbf{Time Allowed: 30 Minutes \hfill Maximum Marks: 25}

\vspace{0.4cm}
\noindent\textbf{Part A: Conceptual Multiple Choice (5 $\times$ 1 = 5 Marks)}
\begin{enumerate}[leftmargin=1.2em, itemsep=2pt]
    \item What does $P = F \times S \times C$ calculate in user feedback analysis?
    \item What is the main difference between Design Refactoring and Product Re-Design?
    \item What do the four letters in the NABC pitching framework represent?
    \item What is the standard duration of an effective elevator pitch?
    \item What is regression testing in the context of design iteration?
\end{enumerate}

\vspace{0.3cm}
\noindent\textbf{Part B: Short Calculation \& Method Questions (2 $\times$ 5 = 10 Marks)}
\begin{enumerate}[leftmargin=1.2em, itemsep=4pt]
    \setcounter{enumi}{5}
    \item Calculate the priority score for two usability bugs:
    \begin{itemize}
        \item Bug 1: A broken checkout button affecting $F=4$ users, Severity $S=5$, Confidence $C=0.9$.
        \item Bug 2: A font misalignment affecting $F=5$ users, Severity $S=1$, Confidence $C=0.8$.
    \end{itemize}
    Explain which bug must be prioritized and why.
    \item Explain the 4 stages of the continuous improvement PDCA cycle.
\end{enumerate}

\vspace{0.3cm}
\noindent\textbf{Part C: Practical Elevator Pitch Formulation (1 $\times$ 10 = 10 Marks)}
\begin{enumerate}[leftmargin=1.2em, itemsep=4pt]
    \setcounter{enumi}{7}
    \item Write a complete, professional 45-second **Elevator Pitch** for an *IoT Smart Pill Dispenser for Alzheimer's Patients* using the Universal Positioning Template. Ensure it highlights Target User, Urgent Need, Category, Core Benefit, Differentiator, Empirical Evidence, and Call-to-Action.
\end{enumerate}
\end{chaptertest}

\begin{solutionbox}[Step-by-Step Unit Test Solutions]
\textbf{Part A Solutions:}
\begin{enumerate}[leftmargin=1.2em, itemsep=2pt]
    \item User feedback Priority Score ($F = \text{Frequency}, S = \text{Severity}, C = \text{Confidence}$).
    \item Refactoring improves internal quality without altering behavior; Re-design overhauls core functional architecture.
    \item Need, Approach, Benefits per costs, Competition/alternatives.
    \item 30 to 60 seconds.
    \item Retesting previously validated features to ensure recent changes have not introduced new defects.
\end{enumerate}

\textbf{Part B Solutions:}
\begin{enumerate}[leftmargin=1.2em, itemsep=3pt]
    \setcounter{enumi}{5}
    \item \textbf{Bug Priority Calculation:}
    \begin{itemize}
        \item Bug 1: $P_1 = 4 \times 5 \times 0.9 = \mathbf{18.0}$.
        \item Bug 2: $P_2 = 5 \times 1 \times 0.8 = \mathbf{4.0}$.
        \item *Decision:* Bug 1 must be prioritized immediately because it is a critical task blocker ($S=5$) that prevents completed purchases, whereas Bug 2 is a minor cosmetic defect ($S=1$).
    \end{itemize}
    \item \textbf{PDCA Cycle:}
    \begin{itemize}
        \item *Plan:* Formulate hypothesis and design the modification.
        \item *Do:* Execute the prototype modification or test pilot.
        \item *Check:* Measure performance against success criteria.
        \item *Act:* Standardize successful changes or iterate further.
    \end{itemize}
\end{enumerate}

\textbf{Part C Solution:}
\begin{enumerate}[leftmargin=1.2em, itemsep=3pt]
    \setcounter{enumi}{7}
    \item \textbf{IoT Smart Pill Dispenser Elevator Pitch:}
    \begin{quote}
    ``For **elderly Alzheimer's patients living independently** who **frequently forget or double-dose their critical daily medications**, our **DoseGuard** is an **IoT-connected automated medication dispenser** that **guarantees exact medication adherence through timed audio-visual chimes and motorized lock-release compartments**. Unlike **traditional static pillboxes and generic smartphone alarms**, DoseGuard **automatically alerts remote family caregivers via SMS if a scheduled dosage is missed after 15 minutes**. In our hospital pilot testing with 40 senior patients, DoseGuard **reduced missed medication incidents from $38\%$ to zero across 60 days**. We are currently seeking **clinical hospital partnerships for our Phase 2 deployment**.''
    \end{quote}
\end{enumerate}
\end{solutionbox}
